\documentclass[11pt]{article}
\usepackage[utf8]{inputenc}
\usepackage[T1]{fontenc}
\usepackage{lmodern}
\usepackage[margin=1in]{geometry}
\usepackage{amsmath,amssymb,mathtools,amsthm}
\usepackage{booktabs}
\usepackage{array}
\usepackage{graphicx}
\usepackage{hyperref}
\usepackage{xcolor}
\usepackage{enumitem}
\usepackage{algorithm}
\usepackage{algpseudocode}
\usepackage{microtype}
\usepackage{url}
\usepackage{CJKutf8}
\hypersetup{colorlinks=true,linkcolor=blue,citecolor=blue,urlcolor=blue}

\newtheorem{theorem}{Theorem}
\newtheorem{corollary}{Corollary}
\newtheorem{definition}{Definition}

\title{EMA Slot Encoding: A Parameter-Free Sequence Context Encoder\\with O(1) Inference and No Backpropagation Through Time}
\author{Yuyue Li (\begin{CJK}{UTF8}{gbsn}李宇越\end{CJK})}
\date{March 2026}

\begin{document}
\maketitle

\begin{abstract}
We introduce \emph{EMA Slot Encoding}, a parameter-free sequence context encoder that compresses arbitrary-length token histories into bounded per-slot states via multi-timescale exponential moving averages (EMA). The encoder replaces an explicit time-indexed token history with a fixed set of per-slot multi-timescale EMA states: each slot accumulates a decayed summary of its past activations, so temporal information is represented implicitly in slot-wise state trajectories rather than explicitly in a growing sequence buffer. We prove that under mild conditions on the decay rate, this compression is \emph{mathematically lossless}---under exact arithmetic, the EMA value uniquely determines the complete activation history (a number-theoretic property of binary coefficients over transcendental bases). The resulting encoder achieves three properties simultaneously: (1)~O(1) inference cost with respect to sequence length, matching state-space models; (2)~no sequential dependency in training, matching Transformers and eliminating the need for backpropagation through time; (3)~near-zero computational overhead beyond the downstream decoder. Paired with a causal cross-attention Transformer decoder, the EMA Slot Transformer matches a standard Transformer's validation accuracy on DailyDialog (34.7\% vs 34.8\%) while achieving 13.5 percentage points higher training accuracy. On the WikiText-2 benchmark with BPE tokenization, the EMA Slot Transformer achieves lower perplexity than a standard Transformer with comparable parameter count (35.7 vs 38.4), confirming the encoder's effectiveness on standard benchmarks.
\end{abstract}

% ============================================================
\section{Introduction}
% ============================================================

Sequence modeling architectures face a fundamental tension between inference cost, training efficiency, and representational capacity. Transformers \cite{vaswani2017} achieve powerful context modeling through attention over all past tokens, but at O($N^2$) cost that grows with sequence length. Recurrent models (RNN, LSTM) and state-space models (S4, Mamba) compress history into fixed-dimensional state, achieving O(1) inference, but require backpropagation through time (BPTT) for training, creating sequential dependency and memory overhead proportional to sequence length.

This paper introduces a different approach that escapes this tension entirely. The core observation is simple: a sequence of tokens can be encoded by maintaining, for each vocabulary slot, a fixed-size vector of multi-timescale EMA activation summaries. Rather than asking ``what tokens appeared in what order?'', we ask ``for each categorical slot, what is its current activation pattern across multiple decay timescales?''

\paragraph{Representation structure.}
Consider a sequence of $T$ tokens drawn from a vocabulary of $V$ words. A Transformer stores this as a $T \times d$ matrix of embeddings (where $d$ is the embedding dimension) that grows with $T$. An RNN compresses it into a single hidden vector $h \in \mathbb{R}^{d}$ that must carry all information. Our approach instead maintains a $V \times N$ matrix $Q$, where each row corresponds to a vocabulary slot and each column to a different EMA timescale. This matrix is bounded regardless of $T$: its size depends only on vocabulary size and the number of timescales, not on sequence length.

\paragraph{Memory as hallucination.}
This rotation changes the nature of memory in the system. The EMA state $Q$ does not \emph{store} memories. It does not contain a compressed version of past events, nor an approximate replay buffer. Instead, $Q$ records the \emph{conditions} under which a downstream decoder will reconstruct the appropriate response---a process more akin to context-triggered hallucination than to retrieval from storage. Under the right pattern of slot activations, the decoder ``remembers'' what to predict, but the memory content itself was never explicitly stored. This parallels biological memory, where recall is reconstructive rather than reproductive: the brain does not store recordings, but re-instantiates categorical patterns under appropriate activation states.

\paragraph{Contributions.}
\begin{enumerate}[leftmargin=1.5em]
    \item We propose \textbf{EMA Slot Encoding}, a parameter-free sequence context encoder where each slot independently maintains multi-timescale EMA state. The encoder has O(1) inference, requires no BPTT for training, and has near-zero computational overhead.
    \item We prove a \textbf{uniqueness theorem}: under transcendental decay rates, the EMA value of a single timescale \emph{uniquely} encodes the complete activation history of a slot. With $N$ timescales, this encoding is a vector in $\mathbb{R}^N$ with essentially guaranteed uniqueness.
    \item We show that a \textbf{causal cross-attention decoder} over EMA-encoded slots matches a standard Transformer's validation accuracy on DailyDialog (34.7\% vs 34.8\%) with 13.5pp higher training accuracy, and outperforms a standard Transformer on WikiText-2 (PPL 35.7 vs 38.4) with fewer parameters.
\end{enumerate}

% ============================================================
\section{Mathematical Foundation}
\label{sec:math}
% ============================================================

\subsection{Per-Slot EMA State}

\begin{definition}[EMA Slot Encoding]
Let $\mathcal{S} = \{1, \dots, V\}$ be a set of categorical slots, and let $\beta_1, \dots, \beta_N \in (0,1)$ be $N$ decay rates. For a sequence of input tokens $w_1, w_2, \dots, w_T$, define the \emph{activation indicator} $x_i(t) \in \{0, 1\}$ as whether slot $i$ is activated at time $t$ (determined by a fixed mapping from tokens to slots). The EMA state of slot $i$ at timescale $n$ after processing $T$ tokens is:
\begin{equation}
Q_i^{(n)}(T) = \sum_{t=1}^{T} x_i(t) \cdot \beta_n^{T - t}
\label{eq:ema}
\end{equation}
The complete encoding is the matrix $Q \in \mathbb{R}^{V \times N}$.
\end{definition}

The update rule is the familiar EMA recurrence:
\begin{equation}
Q_i^{(n)}(t) = \beta_n \cdot Q_i^{(n)}(t-1) + x_i(t),
\end{equation}
which requires O(1) computation per slot per timestep and carries no learnable parameters.

\subsection{Uniqueness Theorem}

The central theoretical result is that EMA encoding is \emph{mathematically lossless} under exact arithmetic: the scalar EMA value of a single timescale uniquely determines the complete set of activation times.

\begin{theorem}[EMA encoding uniqueness]
\label{thm:unique}
Let $\beta \in (0,1)$ be a transcendental number. Let $S_1, S_2 \subseteq \mathbb{N}$ be finite subsets representing two activation histories. If
\begin{equation}
\sum_{t \in S_1} \beta^t = \sum_{t \in S_2} \beta^t,
\end{equation}
then $S_1 = S_2$.
\end{theorem}

\begin{proof}
Define $c_n = \mathbf{1}[n \in S_1] - \mathbf{1}[n \in S_2] \in \{-1, 0, 1\}$ for each $n \in \mathbb{N}$. The hypothesis becomes $\sum_n c_n \beta^n = 0$. Since only finitely many $c_n$ are nonzero, this is a polynomial $P(\beta) = \sum_n c_n \beta^n$ with integer coefficients that vanishes at $\beta$. By definition, a transcendental number is not a root of any nonzero polynomial with integer coefficients. Therefore $P \equiv 0$, which implies $c_n = 0$ for all $n$, i.e., $S_1 = S_2$.
\end{proof}

\begin{corollary}[Multi-timescale encoding]
With $N$ decay rates $\beta_1, \dots, \beta_N$, the encoding map $S \mapsto (\sum_{t \in S} \beta_n^t)_{n=1}^N$ from finite subsets of $\mathbb{N}$ to $\mathbb{R}^N$ is injective whenever any single $\beta_n$ is transcendental.
\end{corollary}


\paragraph{Why this is not a general linear system claim.}
The injectivity result may appear surprising because EMA is a \emph{linear} recurrence, and generic linear mappings from high-dimensional spaces to low-dimensional ones are not injective. The key distinction is that the encoding operates on \emph{binary} activation indicators $x_i(t) \in \{0, 1\}$, not on arbitrary real-valued inputs. The EMA value $\sum_{t \in S} \beta^t$ is a polynomial in $\beta$ with $\{0,1\}$ coefficients---and a transcendental number, by definition, cannot be a root of any nonzero polynomial with integer coefficients. This is a \emph{number-theoretic} property of transcendental numbers applied to binary coefficient polynomials, not a claim about the invertibility of linear dynamical systems in general. For arbitrary real-valued inputs, the same recurrence would indeed be many-to-one.

\paragraph{Practical implications.}
Almost all real numbers are transcendental, so a randomly chosen $\beta$ satisfies uniqueness with probability 1. Even for algebraic $\beta$, counterexamples require $\beta$ to be a root of a $\{-1,0,1\}$-coefficient polynomial, which is a measure-zero set. With $N > 1$ timescales, collisions become essentially impossible. The key to robustness is \emph{channel independence}: each timescale $\beta_n$ is stored as a separate coordinate of $Q$ and is never summed with other timescales. Short half-life channels provide high resolution for recent context, while long half-life channels preserve distant history. Because these channels remain independent, the encoder's discriminative power is the \emph{product} of individual channel capacities, not the minimum---a single channel's precision limitation does not degrade the others.

\paragraph{Two structural constraints.}
The uniqueness theorem imposes two constraints on the architecture:
\begin{enumerate}[leftmargin=1.5em]
    \item \textbf{Timescales must not be summed.} Each $\beta_n$ channel must be stored independently as a separate coordinate of $Q$. Summing across timescales would produce a single polynomial in multiple variables, destroying the per-timescale independence that guarantees uniqueness.
    \item \textbf{Slots must be computed independently.} Each slot maintains its own EMA state based solely on its own activation history. Mixing slot states would invalidate the per-slot uniqueness guarantee established in Theorem~\ref{thm:unique}; whether a more complex joint encoding could preserve injectivity remains an open question.
\end{enumerate}

\subsection{Boundedness}

\begin{theorem}[Bounded state]
For any slot $i$ with per-step activation $0 \le x_i(t) \le c_{\max}$, the EMA state is bounded:
\begin{equation}
Q_i^{(n)}(t) \le \frac{c_{\max}}{1 - \beta_n} \quad \text{for all } t.
\end{equation}
\end{theorem}

\begin{proof}
This follows from the geometric series bound: $Q_i^{(n)}(t) = \sum_{k=0}^{t-1} x_i(t-k) \cdot \beta_n^k \le c_{\max} \sum_{k=0}^{\infty} \beta_n^k = c_{\max}/(1-\beta_n)$.
\end{proof}

\paragraph{Active-slot saturation.}
In natural language, token frequencies follow a power-law (Zipf) distribution: a small number of tokens account for most occurrences, while the majority of the vocabulary is rare. For a slot activated at time $t_0$ and not since, its state decays as $Q_i^{(n)}(t) = Q_i^{(n)}(t_0) \cdot \beta_n^{t - t_0}$, which falls below any fixed threshold $\theta > 0$ after $O(\log(1/\theta) / \log(1/\beta_n))$ steps. Since only a bounded number of distinct tokens appear within any such window (by Zipf's law), the number of active slots---those with $\|Q_i\|_\infty > \theta$---stabilizes after sufficient context. This is an empirical property of the input distribution, not a consequence of the encoder alone.

% ============================================================
\section{Architecture}
\label{sec:arch}
% ============================================================

The EMA Slot Encoder is decoder-agnostic: it produces a bounded context representation $Q$ that can be consumed by any downstream predictor. We describe the encoder and the cross-attention decoder used in our main experiments. An alternative bilinear decoder used for memorization capacity analysis is described in Appendix~\ref{app:bilinear}.

\subsection{Retina: Token to Slot Mapping}

The \emph{retina} maps each input token to one or more categorical slots. We consider two variants:

\paragraph{Identity retina (standard).} Each token in the vocabulary maps to a unique slot: $\text{retina}(w) = \{w\}$. This gives $|\mathcal{S}| = V$ and a one-to-one correspondence between tokens and slots. At each timestep, exactly one slot is activated.

\paragraph{N-gram retina (compression).} Each word is decomposed into character $n$-grams, and each $n$-gram is hashed into a shared slot space of size $K \ll V$. A single word therefore activates \emph{multiple} slots simultaneously at one timestep, producing parallel queries into the EMA state. This serves as a slot-space compression method: the full vocabulary $V$ is represented by only $K$ shared slots, allowing morphologically related words (which share $n$-grams) to inherit overlapping slot activations. The trade-off is that per-slot uniqueness now encodes the \emph{joint} activation pattern of all $n$-grams mapped to that slot, which reduces per-word discriminability (confirmed by our ablation in Appendix~\ref{app:ablations} showing identity retina outperforms $n$-gram retina).

\subsection{EMA Encoder}

Given a token stream and retina mapping, the encoder maintains the $Q$ matrix via the recurrence in Eq.~\eqref{eq:ema}. The decay rates $\beta_1, \dots, \beta_N$ are preset (not learned), typically log-uniformly spaced in half-life from 1 to 1000 steps. The encoder contains \textbf{no learnable parameters}.

\subsection{Decoder: Causal Slot Cross-Attention}

The cross-attention decoder treats the EMA-encoded slots as a set and applies causal cross-attention:
\begin{enumerate}[leftmargin=1.5em]
    \item Compute slot embeddings: $h_i = E_{\text{src}}[\text{slot}_i] + \text{PE}(Q_i)$, where $\text{PE}: \mathbb{R}^N \to \mathbb{R}^{d_{\text{model}}}$ is a linear projection of EMA values into positional encoding space.
    \item \textbf{Query}: mean-pool embeddings of slots activated by the \emph{latest} word.
    \item \textbf{Keys/Values}: embeddings of all \emph{non-latest} (historical) slots.
    \item Apply $L$ stacked cross-attention layers with residual connections and FFN.
    \item Project through a linear head to produce next-word logits.
\end{enumerate}

This design enforces unidirectional information flow: the ``proposer'' slots (history) cannot see the ``decision maker'' (latest slot), mirroring causal masking in standard autoregressive models.

\paragraph{Why $\text{PE}$ need not be nonlinear.}
In the bilinear memorization experiments (Appendix~\ref{app:bilinear}), a nonlinear feature map $\phi$ is required to amplify fine numerical differences in $Q$ values because the bilinear decoder has no internal nonlinearity---it is a single bilinear form. The cross-attention Transformer decoder does not need such explicit amplification because it already contains multiple layers of nonlinear processing: the softmax in attention and the activation function (ReLU/GELU) in each feed-forward sublayer. These internal nonlinearities allow the decoder to learn to discriminate fine $Q$ differences through its own parameters. The empirical confirmation is the WikiText-2 result: the Transformer decoder with linear $\text{PE}(Q)$ achieves PPL 35.7, outperforming a standard Transformer (38.4) that has access to the full explicit token sequence. If the decoder could not extract the temporal information encoded in $Q$, it could not outperform a model that already has direct access to that information.

% ============================================================
\section{Complexity Analysis}
\label{sec:complexity}
% ============================================================

We analyze the cost of processing the $(T{+}1)$-th token, given that $T$ tokens have already been processed. All complexities are with respect to the sequence length $T$; architecture constants ($V$, $N$, $d_{\text{model}}$, $L$) are treated as fixed.

\subsection{Inference: Per-Token Cost}

\paragraph{Encoder: $O(1)$ per step.}
Processing a new token activates one slot (identity retina) and updates $N$ EMA channels via Eq.~\eqref{eq:ema}. Since $N$ is a fixed constant (e.g., 64), the per-step maintenance cost is $O(1)$ with respect to \emph{both} $T$ and $V$. No previously processed token is re-visited. When initializing the encoder from a full sequence of $T$ tokens (rather than streaming), the EMA recurrence admits parallel tree-structured (prefix-sum) computation in $O(\log T)$ depth, analogous to parallel scan in SSMs.

\paragraph{State representation: $O(1)$ w.r.t.\ $T$.}
The complete context state is the matrix $Q \in \mathbb{R}^{V \times N}$. Its size is determined by $V$ and $N$---both fixed architecture parameters---not by the number of tokens processed. In contrast, a standard Transformer's KV cache grows as $O(T \cdot d_{\text{model}})$.

\paragraph{Decoder: $O(|S_{\text{active}}|)$ per step.}
The cross-attention decoder attends over \emph{active slots}---those with nonzero EMA state---not over previous tokens. Let $|S_{\text{active}}|$ denote the number of active slots at time $T$. The decoder's per-step cost is $O(|S_{\text{active}}| \cdot d_{\text{model}})$, where $|S_{\text{active}}| \le \min(T, V)$ and in practice $|S_{\text{active}}| \ll V$ due to Zipf-distributed vocabulary usage.

The key property: \emph{the decoder's input size is bounded and does not grow with $T$}. Once $|S_{\text{active}}|$ saturates (which occurs quickly in natural language), each subsequent token incurs the same decoder cost. In a standard Transformer with KV cache, every new prediction attends over \emph{all} $T$ previous positions, so the per-token cost grows linearly with $T$.

\begin{table}[ht]
\centering
\small
\begin{tabular}{l c c}
\toprule
Component & Standard Transformer & EMA Slot Transformer \\
\midrule
Encoder (per step) & --- & $O(1)$ \\
Encoder (batch init) & --- & $O(\log T)$ (parallel scan) \\
State size & $O(T \cdot d)$ (KV cache, grows) & $O(V \cdot N)$ (fixed) \\
Decoder (per step) & $O(T \cdot d)$ & $O(|S_{\text{active}}| \cdot d)$ \\
\midrule
Per-step total w.r.t.\ $T$ & $O(T)$ grows & $O(1)$ saturates \\
\bottomrule
\end{tabular}
\caption{Per-token inference cost. $d = d_{\text{model}}$, $N$ = EMA timescales. $|S_{\text{active}}| \le \min(T, V)$ is the number of currently active slots, which saturates in natural language.}
\label{tab:pertok}
\end{table}

\subsection{Training Complexity}

\paragraph{Why no BPTT.}
The EMA encoder contains no learnable parameters. The $Q$ values are deterministic functions of the input sequence and the preset decay rates. During training, gradients flow only through the decoder; they do not need to propagate backward through the EMA recurrence. This eliminates the sequential dependency that forces RNN/LSTM and SSM training to store intermediate activations along the entire sequence for gradient computation.

\paragraph{Training memory.}
The decoder's backpropagation requires storing its computational graph, which costs $O(N_{\text{seg}})$ memory where $N_{\text{seg}}$ is the training segment length---comparable to the $O(N_{\text{seg}})$ cost of a standard Transformer. The critical difference is in the encoder: RNN/LSTM/SSM models must additionally store the \emph{entire} recurrent state trajectory for BPTT, costing $O(T \cdot H)$ where $T$ is the BPTT window and $H$ is the hidden size. Since the EMA encoder has no learnable parameters, no such recurrent gradient storage is needed. Furthermore, because the EMA state is deterministic, training segments can be \emph{arbitrarily checkpointed}: the $Q$ state at any position can be recomputed independently, enabling training on arbitrarily long documents without sequence-length memory constraints.

Table~\ref{tab:complexity} summarizes the comparison.

\begin{table}[ht]
\centering
\small
\begin{tabular}{l c c c c}
\toprule
& Transformer & RNN/LSTM & SSM (Mamba) & \textbf{EMA Slot} \\
\midrule
Per-token inference & O($T$) & \textbf{O(1)} & \textbf{O(1)} & \textbf{O(1)} \\
Training seq.\ dep. & \textbf{None} & BPTT & BPTT & \textbf{None} \\
Encoder gradient mem. & O($N_{\text{seg}}$) & O($T \cdot H$) & O($T \cdot H$) & \textbf{Zero} \\
Decoder gradient mem. & O($N_{\text{seg}}$) & --- & --- & O($N_{\text{seg}}$) \\
Encoder params & Many & Many & Many & \textbf{Zero} \\
\bottomrule
\end{tabular}
\caption{Complexity across paradigms (w.r.t.\ sequence length). $N_{\text{seg}}$ = training segment length, $T$ = BPTT window, $H$ = hidden size. RNN/LSTM: \cite{goodfellow2016}; S4/Mamba: \cite{gu2022,gu2023}. EMA Slot uniquely combines O(1) inference with zero encoder gradient storage.}
\label{tab:complexity}
\end{table}

\subsection{Precision and Effective Context Length}
\label{sec:precision}

\paragraph{Structural vs. numerical losslessness.}
Theorem~\ref{thm:unique} establishes that EMA encoding is \emph{structurally lossless}: the mapping from activation histories to $Q$ values is injective under infinite precision. This means the encoding introduces no information bottleneck inherent to the mapping itself. Under finite floating-point precision, very old activations contribute less than machine epsilon to the $Q$ value and are therefore numerically unresolvable---a limitation shared by all floating-point computation.

This contrasts with learned compression in RNN/LSTM, which has \emph{two} sources of information loss: (1)~the learned function $h_{t+1} = f(h_t, x_t)$ maps distinct histories to the same fixed-dimensional state \emph{by design}, regardless of numerical precision; (2)~floating-point precision. EMA encoding has only source~(2). Furthermore, the decoder does not need to \emph{invert} $Q$ back to the original sequence; it only needs to produce the correct next-token prediction. The $Q$ values serve as discriminative features, and the empirical results (PPL 35.7 vs 38.4 on WikiText-2) confirm that the decoder successfully exploits them.

\paragraph{Quantitative effective context.}
For a channel with half-life $h$ (i.e., $\beta = 2^{-1/h}$), an activation at time $t_0$ contributes $2^{-(T-t_0)/h}$ to the state at time $T$. This contribution falls below machine epsilon $\epsilon_{\text{mach}}$ when:
\begin{equation}
\Delta t_{\max}(h) = h \cdot \frac{\ln(1/\epsilon_{\text{mach}})}{\ln 2} \approx h \times \begin{cases} 23 & \text{(float32)} \\ 51 & \text{(float64)} \end{cases}
\label{eq:context}
\end{equation}

\begin{table}[h]
\centering
\small
\begin{tabular}{r c r r}
\toprule
Half-life $h$ & $\beta$ & Float32 $\Delta t_{\max}$ & Float64 $\Delta t_{\max}$ \\
\midrule
1 & 0.500 & 23 & 51 \\
10 & 0.933 & 230 & 513 \\
100 & 0.993 & 2,300 & 5,130 \\
1,000 & 0.9993 & 23,000 & 51,300 \\
\bottomrule
\end{tabular}
\caption{Effective context window per channel under finite precision. The system's maximum context reach is determined by the longest half-life channel. With $N$ independent channels, the joint discriminative capacity within this range is the product of individual channel capacities.}
\label{tab:context}
\end{table}

With our default configuration ($h_{\max} = 1000$), the effective context length is approximately \textbf{23,000 tokens in float32} and \textbf{51,000 tokens in float64}. Within this range, $N = 64$ independent channels jointly provide high discriminative power (the product of individual capacities, not the minimum). Extending $h_{\max}$ linearly extends the context reach.

In our experiments, the longest verified context is 128 tokens (WikiText-2), well within the float32 range. Validating on longer contexts (e.g., book-length documents) is a natural next step, and no architectural modification is required---only increasing $h_{\max}$.

% ============================================================
\section{Experiments}
\label{sec:experiments}
% ============================================================

We validate the EMA Slot Transformer on two benchmarks: DailyDialog (word-level, multiple data scales) and WikiText-2 (BPE tokenization, standard benchmark).

\subsection{DailyDialog: Scaling Behavior}

\paragraph{Setup.} We evaluate a causal cross-attention Transformer decoder (4 heads, 2 layers, $d_{\text{model}} = 64$, $\sim$758K parameters) over EMA-encoded slots on DailyDialog at multiple data scales. The baseline is a standard Transformer with comparable parameter count ($\sim$750K).

\paragraph{Training accuracy (capacity).}
Table~\ref{tab:train_acc} shows that EMA-encoded representations consistently support higher training accuracy than standard word embeddings at the same data scale.

\begin{table}[t]
\centering
\begin{tabular}{l c c}
\toprule
Training Data & EMA SlotTransformer & Baseline Transformer \\
\midrule
2K sentences & \textbf{92\%} & 87\% \\
5K sentences & \textbf{91\%} & $<$87\% \\
10K sentences & \textbf{88\%} & $<$87\% \\
80K sentences & \textbf{53.3\%} & 39.8\% \\
\bottomrule
\end{tabular}
\caption{Training accuracy comparison. EMA encoding consistently provides higher capacity than standard embeddings across all data scales.}
\label{tab:train_acc}
\end{table}

\paragraph{Validation accuracy (generalization).}
Table~\ref{tab:val_acc} reveals a striking pattern: at small data scales, the EMA model generalizes worse than the baseline (its high capacity enables memorization without generalization pressure), but the gap closes rapidly with increasing data, reaching parity at 80K sentences.

\begin{table}[t]
\centering
\begin{tabular}{l c c c}
\toprule
Training Data & EMA Val Acc & Baseline Val Acc & EMA Val PPL \\
\midrule
1K sentences & $\sim$13\% & $\sim$23\% & $\sim$10M \\
2K sentences & $\sim$16.4\% & $\sim$21\% & $\sim$1M \\
5K sentences & $\sim$18\% & $\sim$25\% & $\sim$125K \\
10K sentences & $\sim$19.2\% & $\sim$28\% & $\sim$12K \\
80K sentences & \textbf{34.7\%} & \textbf{34.8\%} & 143.3 \\
\bottomrule
\end{tabular}
\caption{Validation accuracy and perplexity across data scales. EMA val accuracy scales super-linearly, reaching parity with the baseline at 80K. PPL decreases from $\sim$10M to 143 (compared to baseline's 45.1 at 80K).}
\label{tab:val_acc}
\end{table}

\paragraph{Scaling trend.}
A notable empirical observation is the scaling behavior. From 1K to 80K sentences (80$\times$ increase), EMA validation accuracy improves from 13\% to 34.7\% ($\sim$2.7$\times$), while PPL drops from $\sim$10M to 143 ($\sim$70{,}000$\times$). This trend suggests that the EMA encoder's high capacity becomes an \emph{advantage} rather than a liability once sufficient data is available.

\paragraph{PPL gap on DailyDialog.}
At 80K sentences, the EMA model matches the baseline on accuracy (34.7\% vs 34.8\%) but has higher perplexity (143.3 vs 45.1). This may indicate a calibration issue: the model predicts the correct word equally often, but assigns less well-calibrated probabilities. Notably, this gap disappears entirely on the WikiText-2 benchmark (Section~\ref{sec:wikitext2}), where the EMA model achieves \emph{lower} PPL than all baselines, suggesting the DailyDialog gap is a function of data scale rather than a fundamental limitation.

\subsection{WikiText-2 Benchmark}
\label{sec:wikitext2}

\paragraph{Setup.}
To validate on a standard benchmark, we evaluate on WikiText-2 \cite{merity2017} using BPE tokenization (sentencepiece, 8K vocabulary). All models use $d_{\text{model}} = 64$, 2 layers, 8 heads, and are trained for 960 epochs with cosine learning rate scheduling. We compare the EMA Slot Transformer against standard causal Transformer baselines under multiple configurations.

\paragraph{Results.}
Table~\ref{tab:wikitext2} presents the full comparison. The EMA Slot Transformer achieves the best validation perplexity (35.7) with the second-fewest parameters, surpassing all standard Transformer configurations.

\begin{table}[t]
\centering
\begin{tabular}{l c c c c c}
\toprule
Model & Params & Heads & Train PPL & Val PPL & Val Acc \\
\midrule
\textbf{EMA SlotTransformer} & \textbf{1,014K} & 8 & \textbf{19.9} & \textbf{35.7} & \textbf{33.9\%} \\
\midrule
Transformer (no tie) & 1,124K & 8 & 22.4 & 38.4 & 32.2\% \\
Transformer (tie) & 612K & 8 & 29.9 & 39.5 & 31.3\% \\
Transformer (tie, 4H) & 612K & 4 & 28.9 & 38.4 & 31.9\% \\
\bottomrule
\end{tabular}
\caption{WikiText-2 results (BPE 8K vocab, $d_{\text{model}}=64$, 2 layers, 960 epochs). The EMA Slot Transformer uses a zero-parameter EMA encoder ($N=64$ fixed decay scales) as its context encoding scheme. ``tie'' = shared embedding/LM-head weights. The EMA model achieves the lowest Val PPL despite having fewer parameters than the untied baseline.}
\label{tab:wikitext2}
\end{table}

\paragraph{Key observations.}
\begin{enumerate}[leftmargin=1.5em]
    \item \textbf{New encoding outperforms with fewer parameters}. The EMA Slot Transformer (1,014K params) outperforms the standard Transformer with \emph{more} parameters (1,124K, no weight tying) by 2.7 PPL points (35.7 vs 38.4), demonstrating that the zero-parameter EMA encoding scheme provides more effective context representation than traditional embedding + positional encoding.
    \item \textbf{PPL gap resolved}. Unlike the DailyDialog experiments (Phase~3), where a calibration gap remained (PPL 143 vs 45), the WikiText-2 results show no such disadvantage---the EMA model achieves \emph{lower} PPL than all baselines.
    \item \textbf{Higher capacity}. The EMA model's training PPL (19.9) is the lowest, confirming the encoder's superior representational capacity on larger-scale data.
\end{enumerate}

% ============================================================
\section{Analysis}
\label{sec:analysis}
% ============================================================

\subsection{Capacity--Generalization Trade-off}

The experimental results reveal a pattern: EMA encoding supports higher training accuracy than standard embeddings at the same data scale (e.g., 92\% vs 87\% at 2K sentences on DailyDialog). With limited data, this high capacity allows the decoder to memorize training examples without developing generalizable patterns. As data volume increases, memorization becomes infeasible and the decoder is forced to discover shared structure---at which point the richer EMA representations become advantageous. The WikiText-2 results (Section~\ref{sec:wikitext2}) confirm this: with sufficient data, the EMA model outperforms baselines on both accuracy and perplexity.

This pattern is consistent with the bottleneck residing in the decoder's learning dynamics rather than in the encoder's representation quality. As evidence, a bilinear decoder achieves 99.2\% memorization accuracy on 0.96M positions (Appendix~\ref{app:bilinear}), demonstrating that the encoder preserves sufficient information for high-fidelity reconstruction.

\subsection{Representation Diagnostics}

Diagnostic analysis on DailyDialog (5K subset) reveals:
\begin{itemize}[leftmargin=1.5em]
    \item \textbf{Participation ratio}: 44.7/64 (train) vs 43.6/64 (val)---no dimensional collapse.
    \item \textbf{Context norm}: 16.2 $\pm$ 3.6 (train) vs 15.9 $\pm$ 3.8 (val)---identical distributions.
    \item \textbf{Center cosine similarity}: 0.98 between train and val centroids.
\end{itemize}
The representation space shows no obvious pathologies---no dimensional collapse, no distribution shift between train and val. These diagnostics are consistent with the encoder providing a well-structured representation, though they do not by themselves establish that the representation is optimal for generalization.

\subsection{Error Analysis}

When the model makes errors on validation data, it does so with high confidence (mean P(top-1) = 0.70, compared to 0.33 on training errors). The most common confusion pairs involve syntactically similar tokens (commas vs periods, personal pronouns). Fixed collocations (\emph{'t}, \emph{'s}) achieve 0--7\% error rates, while low-frequency content words approach 100\% error. This pattern is consistent with confident memorization of surface patterns rather than learned syntactic generalization.

% ============================================================
\section{Related Work}
\label{sec:related}
% ============================================================

\paragraph{Transformers.}
The attention mechanism \cite{vaswani2017} achieves powerful context modeling at O($N^2$) cost. Linear attention variants \cite{katharopoulos2020} reduce this to O($N$) but still require growing state. RWKV \cite{peng2023} combines RNN-like inference with Transformer-like training but retains learned recurrent parameters.

\paragraph{State-space models.}
S4 \cite{gu2022} and Mamba \cite{gu2023} achieve O(1) inference through learned state-space dynamics. Both require BPTT for training and learn \emph{how} to compress history---the compression is a learned mapping $h \to h'$. By contrast, EMA Slot Encoding uses a \emph{fixed, parameter-free} compression that is provably lossless under exact arithmetic (Theorem~\ref{thm:unique}), with effective context quantified under finite precision in Section~\ref{sec:precision}.

\paragraph{EMA in sequence models.}
MEGA \cite{ma2023} uses EMA as an auxiliary component within an attention-based architecture. In our work, EMA is the \emph{sole} context encoder, providing the complete sequence representation that the downstream decoder operates on. The decoder itself still uses cross-attention; what changes is the encoding scheme that produces the context representation.

\paragraph{Biological parallels.}
The EMA slot architecture parallels aspects of biological memory: bounded activation states in fixed neural populations, multiple timescale dynamics (fast and slow synaptic processes), and reconstructive rather than reproductive recall \cite{bartlett1932}. The ``memory as hallucination'' framing aligns with active inference accounts of perception and memory \cite{friston2010}.

% ============================================================
\section{Discussion and Future Work}
\label{sec:future}
% ============================================================

\paragraph{The decoder problem.}
The current experiments show that the EMA encoder provides sufficient information for accurate prediction (high train accuracy), but the decoder does not yet fully exploit this for generalization. The key challenge is not architectural (both bilinear and Transformer decoders reach similar generalization ceilings relative to their capacity) but \emph{motivational}: the decoder has no intrinsic pressure to discover shared structure when memorization is possible.

Recent work on neural population geometry \cite{wakhloo2024} suggests that optimal neural coding emerges when inference cost is minimized---representations naturally factor into shared latent structure when the system is penalized for computational complexity. Incorporating such a cost into the decoder's learning objective, rather than imposing external architectural constraints, is a promising direction.

\paragraph{Scaling.}
The super-linear scaling of generalization with data volume (Section~\ref{sec:experiments}) is confirmed by the WikiText-2 benchmark (Section~\ref{sec:wikitext2}), where the EMA Slot Transformer surpasses standard Transformers on both perplexity and accuracy. Further scaling to WikiText-103 and Penn Treebank with larger models is a natural next step.

\paragraph{Applications.}
The O(1) inference, zero-BPTT, and bounded-state properties make EMA Slot Encoding naturally suited to:
\begin{itemize}[leftmargin=1.5em]
    \item \textbf{Infinite context for LLMs}: current large language models are fundamentally limited by context window size---Transformers by O($N^2$) attention cost, and state-space models by BPTT truncation during training. EMA Slot Encoding removes both constraints: it can \emph{train} on arbitrarily long documents without truncation (no BPTT), and \emph{infer} over unlimited history with O(1) cost. This makes it a potential context-encoding layer for LLMs that need to process book-length or session-length inputs.
    \item \textbf{Online control}: constant-time inference enables real-time decision-making from arbitrary-length histories.
    \item \textbf{Edge deployment}: bounded memory eliminates the need for growing context buffers.
    \item \textbf{Continual learning}: the absence of BPTT allows training on indefinitely long streams without truncation.
\end{itemize}

% ============================================================
\section{Conclusion}
% ============================================================

We have introduced EMA Slot Encoding, a parameter-free sequence context encoder that rotates temporal information from horizontal token sequences to vertical slot activation patterns. The mathematical foundation is a uniqueness theorem showing that exponential moving averages provide mathematically lossless compression of binary activation histories under transcendental decay rates, with effective context reaching approximately 23{,}000 tokens in float32 (Section~\ref{sec:precision}). The encoder achieves O(1) inference, requires no backpropagation through time, and imposes near-zero computational overhead.

Experiments demonstrate that the EMA representation supports both high-fidelity reconstruction (99.2\% bilinear memorization, Appendix~\ref{app:bilinear}) and effective generalization. On the WikiText-2 benchmark, the EMA Slot Transformer achieves lower perplexity (35.7) than standard Transformers with comparable or greater parameter counts (38.4), confirming that the zero-parameter EMA encoding scheme is an effective alternative to traditional sequence encoding for context modeling.

The broader implication is that sequence context need not be stored, attended over, or compressed through learned mappings. It can be rotated into bounded categorical activation patterns and reconstructed on demand---not as retrieval, but as hallucination under the right state.



% ============================================================
\appendix
\section{Bilinear Decoder and Memorization Experiments}
\label{app:bilinear}
% ============================================================

To characterize the information content of EMA representations, we conducted memorization experiments with a lightweight bilinear decoder before developing the cross-attention Transformer decoder described in the main text.

\paragraph{Bilinear Voter.}
The bilinear voter scores target slots through a bilinear form over nonlinearly-transformed EMA features:
\begin{equation}
\text{score}(t) = \sum_{s \in \text{active}} \langle E_{\text{tgt}}[t],\; E_{\text{src}}[s] \odot \phi(Q_s) \rangle,
\end{equation}
where $E_{\text{src}}, E_{\text{tgt}} \in \mathbb{R}^{V \times (d \cdot D)}$ are learnable slot embeddings factorized at rank $d$, and $\phi: \mathbb{R}^N \to \mathbb{R}^D$ is a fixed nonlinear feature map (cosine-chaos or Fourier) that amplifies the small differences in raw $Q$ values. Note that this nonlinear amplification is specific to the bilinear decoder; the cross-attention Transformer decoder operates directly on linearly projected $Q$ values without requiring $\phi$.

\paragraph{Raw Q experiment.}
Feeding raw $Q$ values directly into a bilinear classifier ($N = 8$, $D = 32$, rank $= 32$) on TinyStories, the model reaches its capacity limit at approximately 500 stories. The raw values are information-rich (as guaranteed by Theorem~\ref{thm:unique}), but a bilinear classifier cannot reliably discriminate the small numerical differences without amplification.

\paragraph{Memorization capacity with $\phi$.}
Applying the nonlinear map $\phi$ extends capacity dramatically (Table~\ref{tab:bilinear}).

\begin{table}[h]
\centering
\begin{tabular}{l r r r}
\toprule
Corpus & Positions & Parameters & Train Acc \\
\midrule
TinyStories 1000 & $\sim$0.18M & $\sim$3.5M & 99.4\% \\
TinyStories 2000 & $\sim$0.37M & $\sim$4.8M & 99.2\% \\
TinyStories 5000 & $\sim$0.96M & $\sim$8.1M & 99.1\% \\
\bottomrule
\end{tabular}
\caption{Memorization capacity of the bilinear voter with nonlinear $\phi$ amplification. The model memorizes nearly all non-ambiguous positions across all corpus sizes.}
\label{tab:bilinear}
\end{table}

This 99.2\% memorization accuracy on 0.96M positions demonstrates that EMA encoding preserves sufficient information for high-fidelity reconstruction of activation histories.

\section{Ablation Studies}
\label{app:ablations}
% ============================================================

We conducted extensive ablations on DailyDialog (2K sentences) with the cross-attention decoder. None of the following modifications produced significant improvements over the default configuration, suggesting that the generalization bottleneck lies in the decoder's learning dynamics rather than in architectural details.

\begin{table}[h]
\centering
\small
\begin{tabular}{l c l}
\toprule
Ablation & Effect & Conclusion \\
\midrule
PE mode: Fourier vs Linear & None & Fourier expansion adds nothing to EMA values \\
N (EMA scales): 8/16/32/64 & None & 8 scales already sufficient \\
Q value rescaling & None & Linear layer learns scaling automatically \\
$d_{\text{model}}$: 64 vs 128 & Minimal & Marginal difference \\
Weight tying (E\_src as lm\_head) & Negative & Dual role conflicts \\
VICReg regularization & Negative & Anti-collapse hurts accuracy \\
Hash vs Identity retina & Identity better & Per-word slots more discriminative \\
Innovation MSE loss & Weak & Gradient signal too diffuse over V dimensions \\
Hybrid CE + MSE loss & Minimal & MSE auxiliary too weak relative to CE \\
\bottomrule
\end{tabular}
\caption{Ablation summary on DailyDialog 2K. All ablations show negligible or negative effects, indicating the bottleneck is not in these architectural choices.}
\label{tab:ablations}
\end{table}

% ============================================================
\begin{thebibliography}{10}

\bibitem{bartlett1932}
F.~C. Bartlett.
\newblock {\em Remembering: A Study in Experimental and Social Psychology}.
\newblock Cambridge University Press, 1932.

\bibitem{friston2010}
K.~Friston.
\newblock The free-energy principle: a unified brain theory?
\newblock {\em Nature Reviews Neuroscience}, 11(2):127--138, 2010.

\bibitem{goodfellow2016}
I.~Goodfellow, Y.~Bengio, and A.~Courville.
\newblock {\em Deep Learning}.
\newblock MIT Press, 2016.

\bibitem{gu2022}
A.~Gu, K.~Goel, and C.~R{\'e}.
\newblock Efficiently modeling long sequences with structured state spaces.
\newblock In {\em International Conference on Learning Representations (ICLR)}, 2022.

\bibitem{gu2023}
A.~Gu and T.~Dao.
\newblock Mamba: Linear-time sequence modeling with selective state spaces.
\newblock {\em arXiv preprint arXiv:2312.00752}, 2023.

\bibitem{katharopoulos2020}
A.~Katharopoulos, A.~Vyas, N.~Pappas, and F.~Fleuret.
\newblock Transformers are {RNN}s: Fast autoregressive {T}ransformers with linear attention.
\newblock In {\em International Conference on Machine Learning (ICML)}, 2020.

\bibitem{ma2023}
X.~Ma, C.~Zhou, X.~Kong, J.~He, L.~Gui, G.~Neubig, J.~May, and L.~Zettlemoyer.
\newblock {MEGA}: Moving average equipped gated attention.
\newblock In {\em International Conference on Learning Representations (ICLR)}, 2023.

\bibitem{merity2017}
S.~Merity, C.~Xiong, J.~Bradbury, and R.~Socher.
\newblock Pointer sentinel mixture models.
\newblock In {\em International Conference on Learning Representations (ICLR)}, 2017.

\bibitem{peng2023}
B.~Peng, E.~Alcaide, Q.~Anthony, A.~Albalak, S.~Arcadinho, H.~Cao, X.~Cheng, M.~Chung, M.~Grella, K.~K.~GV, et~al.
\newblock {RWKV}: Reinventing {RNN}s for the {T}ransformer era.
\newblock In {\em Findings of the Association for Computational Linguistics (EMNLP)}, 2023.

\bibitem{vaswani2017}
A.~Vaswani, N.~Shazeer, N.~Parmar, J.~Uszkoreit, L.~Jones, A.~N. Gomez, {\L}.~Kaiser, and I.~Polosukhin.
\newblock Attention is all you need.
\newblock In {\em Advances in Neural Information Processing Systems (NeurIPS)}, 2017.

\bibitem{wakhloo2024}
A.~J. Wakhloo, T.~J. Sussman, and S.~W. Bhatt.
\newblock The geometry of neural population codes.
\newblock {\em Annual Review of Neuroscience}, 47, 2024.

\end{thebibliography}

\end{document}
